\documentclass[11pt]{article}

\usepackage[margin=1in]{geometry}
\usepackage{amsmath,amssymb,amsthm,mathtools}
\usepackage{booktabs,tabularx,array,longtable}
\usepackage[section]{placeins}
\usepackage{float}
\usepackage{microtype}
\usepackage{graphicx}
\usepackage[numbers,sort&compress]{natbib}
\usepackage[colorlinks=true,allcolors=blue]{hyperref}
\usepackage[nameinlink]{cleveref}
\usepackage[T1]{fontenc}
\usepackage[utf8]{inputenc}
\usepackage{xcolor}
\usepackage{listings}
\usepackage{listingsutf8}
\usepackage{authblk}
\setlength{\emergencystretch}{3em}

% ArXiv-friendly Lean code: no shell escape, minted, or external highlighter.
\definecolor{leankeyword}{HTML}{243B6B}
\definecolor{leantactic}{HTML}{70477A}
\definecolor{leantype}{HTML}{1F6473}
\definecolor{leancomment}{HTML}{39734D}
\definecolor{leanstring}{HTML}{8A4B2A}
\definecolor{leannumber}{HTML}{777777}
\definecolor{leanbackground}{HTML}{F7F7F3}
\lstdefinelanguage{lean}{%
  sensitive=true,
  morecomment=[l]{--},
  morecomment=[s]{/-}{-/},
  morestring=[b]",
  morekeywords=[1]{import,open,namespace,section,end,noncomputable,variable,variables,
    universe,universes,set_option,attribute,axiom,opaque,theorem,lemma,def,abbrev,
    example,instance,structure,class,inductive,notation,local,deriving,extends,
    by,where,let,have,show,fun,match,with,if,then,else,do,return,in,forall},
  morekeywords=[2]{simp,simpa,rw,erw,dsimp,unfold,ring,ring_nf,nlinarith,linarith,
    positivity,norm_num,field_simp,omega,exact,refine,apply,assumption,intro,intros,
    ext,funext,cases,rcases,obtain,constructor,left,right,classical,by_contra,
    contrapose,convert,using,suffices,calc},
  morekeywords=[3]{Prop,Type,Nat,Int,Real,Fin,Set,Submodule,Matrix,EuclideanSpace,
    LinearMap,ContinuousLinearMap,LinearPMap,RCLike,IsTrialResidual,
    IsExactSpectralDecomposition,FormBoundedSylvesterGap,
    NormalizedSymmetricOperatorIdealFamily,SourcePopulationGap},
}
\lstdefinestyle{leanpaper}{%
  language=lean,
  columns=fullflexible,
  backgroundcolor=\color{leanbackground},
  basicstyle=\ttfamily\scriptsize,
  commentstyle=\color{leancomment}\itshape,
  keywordstyle=[1]\color{leankeyword}\bfseries,
  keywordstyle=[2]\color{leantactic},
  keywordstyle=[3]\color{leantype},
  stringstyle=\color{leanstring},
  numberstyle=\tiny\color{leannumber},
  breakatwhitespace=false,
  breaklines=true,
  captionpos=b,
  keepspaces=true,
  numbers=left,
  numbersep=6pt,
  showspaces=false,
  showstringspaces=false,
  showtabs=false,
  tabsize=2,
  aboveskip=0.8em,
  belowskip=0.8em,
  literate=
    {\\KField}{{$\Bbbk$}}1
}
\crefname{lstlisting}{formalization}{formalizations}
\Crefname{lstlisting}{Formalization}{Formalizations}
\renewcommand{\lstlistingname}{Formalization}

\newtheorem{theorem}{Theorem}
\newtheorem{proposition}[theorem]{Proposition}
\newtheorem{remark}[theorem]{Remark}

% Generated by scripts/build_source_census.py; do not edit by hand.
\newcommand{\YWSCensusEntryCount}{24}
\newcommand{\YWSCensusPaperSourceCount}{22}
\newcommand{\YWSCensusExplicitAdditionCount}{2}
\newcommand{\YWSCensusVerifiedCount}{24}
\newcommand{\YWSCensusCorrectedCount}{7}

% Generated by scripts/fetch_openalex_bibliometrics.py; do not edit by hand.
\newcommand{\OpenAlexRetrievedAtUTC}{2026-08-17T20:33:55Z}
\newcommand{\OpenAlexRetrievedDate}{17 August 2026}
\newcommand{\OpenAlexCurrentYear}{2026}
\newcommand{\OpenAlexLatestCompleteYear}{2025}
\newcommand{\OpenAlexDavisKahanCitations}{1228}
\newcommand{\OpenAlexDavisKahanWorkID}{W1970377488}
\newcommand{\OpenAlexDavisKahanUpdatedDate}{2026-08-16T07:02:28.622633}
\newcommand{\OpenAlexDavisKahanCurrentYearCitations}{57}
\newcommand{\OpenAlexDavisKahanLatestCompleteYearCitations}{60}
\newcommand{\OpenAlexYuWangSamworthCitations}{462}
\newcommand{\OpenAlexYuWangSamworthWorkID}{W2088911135}
\newcommand{\OpenAlexYuWangSamworthUpdatedDate}{2026-08-16T07:02:28.622633}
\newcommand{\OpenAlexYuWangSamworthCurrentYearCitations}{22}
\newcommand{\OpenAlexYuWangSamworthLatestCompleteYearCitations}{49}


% Manually populated from Google Scholar on the same date as the OpenAlex observation.
\newcommand{\GoogleScholarDavisKahanCitations}{2086}
\newcommand{\GoogleScholarYuWangSamworthCitations}{1067}

\newcommand{\Lean}{\textsc{Lean}}
\newcommand{\YWS}{Yu--Wang--Samworth}
\newcommand{\DK}{Davis--Kahan}
\newcommand{\TauCeti}{Tau Ceti}
\newcommand{\DKResultCount}{29}
\newcommand{\DKProvedCount}{28}
\newcommand{\DKRefutedCount}{1}
\newcommand{\DKSemanticAcceptedCount}{29}

\title{Formalizing Davis--Kahan Perturbation Theory in Lean}
% Additional authors are expected; keep the confirmed ordering below stable.

\author[1]{Jonathan Crall}
\author[2]{Edward Wang}
\author[1]{Brian Hu}
\author[2]{Carey E. Priebe}

\affil[1]{%
  Kitware, Inc.\\
  \texttt{jon.crall@kitware.com},
  \texttt{brian.hu@kitware.com}
}

\affil[2]{%
  Johns Hopkins University\\
  \texttt{ewang43@jhu.edu},
  \texttt{cep@jhu.edu}
}

\date{Draft -- September 2026}

\begin{document}
\maketitle

\begin{abstract}
We formalize in Lean~4 the operator perturbation theory developed by Davis and
Kahan in their 1970 paper on rotations of eigenspaces.  The development covers
the paper's result-level mathematical content, including the real and complex,
infinite-dimensional, unbounded-operator, and unitarily invariant norm settings
used in the source.  Of the 29 results selected by an explicit source-level
counting rule, 28 are proved at the stated scope.  Proposition~4.4 is false as
printed: we give a machine-checked counterexample in $\mathbb{R}^4$, locate the
invalid inequality in the published proof, and prove a corrected theorem for
$Q$-norms.

The single-angle sine theorem illustrates the analytic scope of the
formalization.  Its formal statement permits unbounded self-adjoint operators,
finite and half-infinite spectral separation, and where-defined unitarily
invariant norms.  Relating this statement to Davis and Kahan's notation also
requires a mathematical identification: the rectangular operator
$(I-F_0F_0^*)E_0$ used in Lean has modulus equal to the positive
$\sin\Theta_0$ operator of the paper, so the two have the same supported
unitarily invariant norms.

Reaching this scope required formal theory for principal angles and direct
rotations, approximation numbers and symmetric operator ideals, Borel spectral
calculus, unbounded self-adjoint operators, real/complex transfer, and Sylvester
equations.  We also formalize the theorem-level claims of Yu, Wang, and
Samworth as a finite-dimensional statistical specialization and verify two
printed defects in that paper.
\end{abstract}

\section{Introduction}
\label{sec:introduction}

Davis and Kahan's perturbation theory for spectral subspaces connects
perturbations of self-adjoint operators to geometric changes in their spectral
subspaces \citep{davis1970rotation,kato1976perturbation,stewart1990matrix}.
The sine-theta inequality is its best-known descendant, especially in matrix
statistics, but the 1970 paper has a broader operator-theoretic scope.  It
develops canonical direct rotations between subspaces, extremal properties
under unitarily invariant norms, Sylvester-operator estimates, generalized sine
and tangent bounds, and results for unbounded self-adjoint operators.  These
ideas remain part of the standard perturbation toolkit
\citep{wedin1972perturbation,chen2021spectral}.  On
\OpenAlexRetrievedDate{}, OpenAlex indexed \OpenAlexDavisKahanCitations{}
citing works for Davis--Kahan and \OpenAlexYuWangSamworthCitations{} for
Yu--Wang--Samworth (YWS); Google Scholar reported
\GoogleScholarDavisKahanCitations{} and \GoogleScholarYuWangSamworthCitations{}
citations, respectively.

Lean~4 is an interactive theorem prover whose kernel checks elaborated proof
terms, and Mathlib is its community mathematical library
\citep{demoura2021lean4,mathlib2020}.  Formalizing Davis--Kahan therefore asks
two mathematical questions.  First, can the stated result be proved from a
precise collection of encoded hypotheses?  Second, does that encoded statement
have the same scope as the result in the paper?  The second question is
substantive here because boundedness, scalar field, spectral-gap geometry,
operator domains, and the class of admissible norms all change the meaning of a
perturbation theorem.

The distinction appears already in the sine theorem.  A finite
interval/exterior gap gives a useful Davis--Kahan estimate, but the source also
proves ordered half-infinite separation cases.  Likewise, replacing an
unbounded operator by a bounded matrix statement preserves the familiar shape
of the inequality while losing part of the theorem.  Our formal endpoint keeps
these cases and makes the domain hypotheses explicit.  A different issue arises
in Section~4: Proposition~4.4 has no proof at its printed scope because the
claim itself is false.

The formalization consequently required operator theory beyond the matrix-level
statement usually associated with Davis--Kahan.  We developed the needed
principal-angle geometry, symmetric operator ideals, spectral calculus,
unbounded-operator interfaces, direct rotations, and Sylvester estimates as
reusable mathematics.  Some of the spectral infrastructure was adapted or
generalized from the Spectra Lean development \citep{bornemann2026spectra}; the
resulting reusable declarations are organized under the \texttt{TauCeti}
namespaces.

YWS \citep{yu2015davis} provides a smaller statistical specialization.  Its
population-gap bound packages the subspace perturbation estimate in a form that
can be combined directly with stochastic control of matrix-estimation error.
We use it as a secondary test of the finite-dimensional spectral geometry and
of source-level checking; its formalization exposed one algebraic typo and one
boundary-convention error.

\paragraph{Contributions.}
\begin{enumerate}
  \item We give a Lean treatment of all \DKResultCount{} result-level targets
        in Davis and Kahan's 1970 paper under an explicit counting rule.
        \DKProvedCount{} are proved at the scope stated in the source, and
        Proposition~4.4 is refuted by a checked counterexample.
  \item We formalize the single-angle sine theorem in its unbounded,
        real/complex, finite-or-half-infinite gap setting with where-defined
        unitarily invariant norms, and prove the modulus identity that connects
        the Lean operator to Davis and Kahan's positive $\sin\Theta_0$ operator.
  \item We formalize the direct-rotation, operator-ideal, Sylvester, generalized
        sine/tangent, and branch-selection theory needed for the rest of the
        paper, building reusable operator-theoretic foundations where the
        required interfaces were absent.
  \item We identify the failure of Davis--Kahan Proposition~4.4 inside its
        advertised short-angle regime and prove a corrected full-displacement
        theorem for $Q$-norms.  As a secondary application, we formalize YWS
        Theorems~1--3 and verify two printed source defects.
\end{enumerate}

\section{Scope of the Davis--Kahan formalization}
\label{sec:coverage}

The formalization covers the result-level mathematical content established in
Davis and Kahan's paper.  The counting rule includes the four unnumbered
Section~2 theorems and every named theorem, proposition, lemma, and corollary
proved by Davis and Kahan; definitions, examples, open questions, proof-local
claims, and results attributed to other sources are excluded.  This gives
\DKResultCount{} targets.  \DKProvedCount{} have checked Lean proofs at the
scope stated in the paper.  Proposition~4.4 is the remaining target and is
refuted as printed.  Appendix~\ref{app:dk-inventory} gives the complete
result-by-result accounting and Appendix~\ref{app:source-review} describes the
statement-review criterion.

The breadth of the source is more informative than the count.  Sections~2 and~6
include sine and tangent perturbation bounds under several forms of spectral
separation.  Section~3 develops the geometry of two subspaces and their direct
rotation.  Section~4 studies extremality under unitarily invariant norms.
Section~5 supplies Sylvester estimates in bounded and unbounded settings, and
Section~8 selects the appropriate spectral branch under smallness hypotheses.
The formalization follows these dependencies through infinite-dimensional
Hilbert spaces where they occur in the source, including the approximation-number
machinery needed to state unitarily invariant norm comparisons without reducing
them to finite matrices.

\section{A full-scope sine-theta formalization}
\label{sec:sine-theta}

\subsection{Trial and exact subspaces}

Davis and Kahan distinguish a reference operator $A$ from a perturbed operator
$A+H$.  For the formal statement we write $A$ for the perturbed full-space
operator, corresponding to the source's $A+H$.

Let $E$ be the ambient Hilbert space.  Let $A_0$ be a self-adjoint operator on a
trial-coordinate space $F$, and let
\[
  E_0:F\to E
\]
be an isometric coordinate map for the trial subspace.  The residual compares
applying the ambient operator after embedding with applying the trial operator
before embedding:
\begin{equation}
  R = AE_0-E_0A_0.
  \label{eq:residual}
\end{equation}
When the operators are unbounded, \cref{eq:residual} is asserted on
$\operatorname{dom}(A_0)$ and $R$ is the bounded residual that agrees there.
The Lean predicate \texttt{IsTrialResidual} packages the isometry, the required
domain inclusion, and this residual identity.

Let $F_0:H\to E$ be isometric coordinates for the desired exact spectral
subspace and let $F_1:G\to E$ be isometric coordinates for its orthogonal
complement.  A self-adjoint operator $\Lambda_1$ on $G$ describes the
complementary spectral block through
\[
  AF_1 = F_1\Lambda_1
\]
on its domain.  The predicate \texttt{IsExactSpectralDecomposition} also records
orthogonality and completeness of the two coordinate ranges.

\subsection{The sine operator and its rectangular representative}

The source angle operator acts on the trial-coordinate space.  If
$P=F_0F_0^*$ is the orthogonal projection onto the desired exact subspace, then
\[
  \cos\Theta_0 = (E_0^*PE_0)^{1/2},
  \qquad
  \sin\Theta_0 = (I-E_0^*PE_0)^{1/2}.
\]
The Lean theorem places the following rectangular map directly inside the
operator ideal:
\begin{equation}
  S=(I-P)E_0.
  \label{eq:rectangular-sine}
\end{equation}
The two operators have different codomains:
$\sin\Theta_0:F\to F$ and $S:F\to E$.  Their relationship is
\[
  S^*S
  = E_0^*(I-P)E_0
  = I-E_0^*PE_0,
\]
so
\begin{equation}
  |S|=(S^*S)^{1/2}=\sin\Theta_0.
  \label{eq:modulus-sine}
\end{equation}
A polar decomposition $S=U|S|$ and the ideal laws imply that $S$ and $|S|$
have the same singular-value data and the same value under every supported
unitarily invariant norm.  Consequently
\[
  N(S)=N(\sin\Theta_0)
\]
whenever the norm is defined.  This modulus identity is the semantic bridge
between the source operator and the Lean conclusion.

\subsection{Gap geometry and the inequality}

The source includes three forms of spectral separation.  The familiar case
places one diagonal block in a finite interval $[\beta,\alpha]$ and the other
outside its open $\delta$-neighborhood.  The Appendix to Section~6 also permits
ordered half-infinite separation: one block may be semibounded above while the
other is semibounded below at distance at least $\delta$, in either order.  The
Lean predicate \texttt{FormBoundedSylvesterGap} has constructors for all three
configurations.  The source spectral hypotheses imply the corresponding
form-bounded separation used by the Sylvester estimate.

The source also writes a common-domain equality for the unbounded trial
configuration.  That equality discharges the forward domain condition consumed
by \texttt{IsTrialResidual}.  Source correspondence is checked from the combined
hypotheses, with this equality supplying the required domain condition.

With these definitions, the formal theorem states the
where-defined inequality
\begin{equation}
  \delta\,N(S)\leq N(R),
  \qquad S=(I-F_0F_0^*)E_0,
  \label{eq:dk-sine}
\end{equation}
for real or complex separable Hilbert spaces and possibly unbounded
self-adjoint $A$, $A_0$, and $\Lambda_1$.  ``Where-defined'' means that the
operator-ideal gauge is finite on both displayed operators.  This matches the
source convention and leaves membership transfer for arbitrary symmetric
operator ideals as a separate theorem.

Appendix~\ref{app:lean-guide} summarizes the general Lean notation used in the
displayed signatures.

\begin{lstlisting}[style=leanpaper,
  caption={Lean signature for the full sine-theta theorem, abbreviated for
  readability.  Here \texttt{S} denotes the inlined rectangular map
  $(I-F_0F_0^*)E_0$; the exact Lean source uses the same theorem name and
  expands \texttt{S} in both membership and gauge expressions.},
  label={lst:dk-sintheta-headline}]
theorem sinTheta_unbounded_formGap_whereDefinedUIN_rclike
    [TopologicalSpace.SeparableSpace E]
    (N : NormalizedSymmetricOperatorIdealFamily K)
    (A : LinearPMap K E E)
    (A0 : LinearPMap K F F)
    (Lambda1 : LinearPMap K G G)
    (E0 : ContinuousLinearMap K F E)
    (F0 : ContinuousLinearMap K H E)
    (F1 : ContinuousLinearMap K G E)
    (R : ContinuousLinearMap K F E)
    (hA : IsSelfAdjoint A)
    (hA0 : IsSelfAdjoint A0)
    (hLambda1 : IsSelfAdjoint Lambda1)
    (htrial : IsTrialResidual A A0 E0 R)
    (hexact : IsExactSpectralDecomposition A Lambda1 F0 F1)
    (delta : Real) (hdelta : 0 < delta)
    (hgap : FormBoundedSylvesterGap A0 Lambda1 delta) :
    N.Mem S ->
    N.Mem R ->
      delta * N.gaugeReal S <= N.gaugeReal R := by
\end{lstlisting}

The theorem contains all three separation configurations encoded by its gap
hypothesis.  The ordered half-infinite cases have distinct hypotheses from
the finite interval/exterior specialization, so source correspondence is checked
against the complete family of gap hypotheses.

\section{Mathematical architecture of the Davis--Kahan development}
\label{sec:dk-rest}

The single-angle estimate depends on mathematical structures that recur through
the 1970 paper.  Formalizing the remaining results made these dependencies
explicit and led to three substantial layers: the geometry of two subspaces,
symmetric operator ideals and majorization, and Sylvester theory for spectral
separation.  The reusable parts of these layers are kept below the
paper-specific theorem package.

\subsection{Direct rotations and the geometry of two subspaces}

Let $P$ and $Q$ be the orthogonal projections onto closed subspaces $U$ and
$V$.  Davis and Kahan organize their geometry through the operators obtained
from $P$, $Q$, and their complements.  In the acute case, the compression
$PQP$ has no kernel on $U$ and determines the cosine of the angle operator;
the complementary compression determines its sine.  Polar decomposition then
supplies the partial isometries that pair the corresponding principal
subspaces.  Combining these pieces gives the direct rotation, an orthogonal or
unitary operator $R$ satisfying $R(U)=V$ and rotating along the principal
planes.

This construction explains several results that look separate when stated only
at theorem level.  Existence and uniqueness in the acute case, the
principal-square-root characterization, commutation with the angle operators,
and reversal symmetry all express compatibility among the same projection and
polar-decomposition data.  The nonacute theory records the additional kernel
conditions needed when principal angles reach $\pi/2$.  For compact angle
operators, the resulting decomposition reduces to the familiar sequence of
principal angles; the general Hilbert-space statements do not require such a
discrete spectrum.

Section~8 returns to this geometry from the spectral side.  Perturbation
smallness and spectral repulsion rule out the wrong branch and place the two
spectral subspaces in the acute regime where the direct rotation is available.
In the formal development this branch-selection argument uses the same
unbounded spectral calculus as the perturbation theorems.  Acuteness therefore
follows from the spectral hypotheses, with no separate matrix-level assumption.

\subsection{Approximation numbers, majorization, and unitarily invariant norms}

Davis and Kahan formulate many extremal statements simultaneously for broad
classes of unitarily invariant norms.  In finite dimensions such a norm can be
described by a symmetric gauge function of the singular values.  The
infinite-dimensional formalization uses approximation numbers as the singular
value sequence of a compact operator and symmetric norming functions to encode
the gauge.  Ky Fan partial sums then provide the bridge from pointwise or
majorization statements to norm inequalities.

This layer is essential in Section~4.  If $R$ is the direct rotation and $W$
is any unitary carrying $U$ onto $V$, the restricted displacement has
singular-value extremality strong enough to imply
unitarily invariant norm minimality.  Davis and Kahan also prove a
majorization theorem for the squared full displacement
$(I-W)^*(I-W)$.  Squaring changes the relevant singular values from
$s_j(I-W)$ to $s_j(I-W)^2$, so a majorization statement at the squared level
does not automatically imply the corresponding statement for the unsquared
trace norm.  Proposition~4.4 attempts precisely that stronger passage.  The
counterexample in Section~\ref{sec:prop44} shows where it fails, while the
$Q$-norm repair identifies a norm class for which the squared theorem does
transfer back to $I-W$.

The same operator-ideal machinery is used in the perturbation inequalities.
For the sine theorem, the statement $\delta N(S)\leq N(R)$ is meaningful for
any supported symmetric operator ideal once both $S$ and the residual belong to
that ideal.  The formal theorem therefore keeps the source's ``where-defined''
convention and does not build an accidental finite-dimensional compactness
assumption into the statement.

\subsection{Sylvester separation and unbounded operators}

The analytic engine behind the sine and tangent estimates is the Sylvester map
\[
  X \longmapsto AX-XB.
\]
When the spectra of self-adjoint $A$ and $B$ are separated by a gap $\delta$,
Davis and Kahan obtain lower bounds of the form
\[
  \delta\,N(X) \leq N(AX-XB).
\]
The perturbation theorems result by identifying the appropriate off-diagonal or
trial-subspace operator $X$ and rewriting its Sylvester residual in terms of the
original perturbation.

The unbounded case forces domain information into this argument.  Expressions
such as $AX-XB$ are only meaningful on the appropriate domains, and the
spectral separation may be encoded by a finite interval/exterior configuration
or by an ordered pair of half-lines.  The reusable development represents
unbounded operators as domain-carrying linear partial maps, develops the
self-adjoint and semibounded interfaces needed by the proof, and connects
spectral hypotheses to a form-bounded Sylvester gap.  Strong spectral cutoffs
reduce parts of the argument to bounded operators; convergence of the relevant
approximation numbers then permits passage back to the unbounded setting.

Section~6 reuses this engine after weakening the trial-coordinate assumptions.
Reflection pinching and direct-sum norm comparisons control the resulting block
operators, while the Sylvester estimate supplies the spectral-gap factor.  This
produces the generalized sine and tangent theorems and the pairwise-gap
square-norm estimate at the scope stated in the source.

\subsection{Formal foundations and ancestry}

Several of the preceding objects had no sufficiently general interface in the
available Lean libraries when this work began.  The development therefore adds
a focused theory of principal angles and projections, approximation numbers and
symmetric operator ideals, polar decomposition and direct rotations, Borel
spectral calculus, unbounded self-adjoint partial maps, and real/complex
transfer.  Complexification is used to prove analytic results once when the
complex spectral theorem is the natural tool and then descend the conclusions
to real Hilbert spaces where Davis and Kahan state them as well.

Mathlib supplies the general algebraic, topological, Hilbert-space, and
continuous-functional-calculus foundation \citep{mathlib2020,dedecker2025cfc}.
Parts of the projection-valued-measure, Borel-calculus, and self-adjoint
infrastructure were copied, adapted, or generalized from Spectra
\citep{bornemann2026spectra}.  Reusable declarations from this project are
organized under the \texttt{TauCeti.*} namespaces and are being prepared for
integration with Tau Ceti \citep{tauceti2026}.  The paper-specific
Davis--Kahan package contains the source statements, examples, and the
Proposition~4.4 counterexample; Appendix~\ref{app:source-review} records the
statement-level review procedure.

\section{Proposition 4.4: refutation and repair}
\label{sec:prop44}

Davis--Kahan Proposition~4.4 asserts, in the real finite-dimensional case, that
when the largest principal angle is at most $\pi/3$, the direct rotation
minimizes every unitarily invariant norm of the full displacement $I-W$ among
unitaries carrying one subspace onto the other.  The statement is false.

\subsection{A four-dimensional counterexample}

Let $E=\mathbb{R}^4$ and
\[
  U=\operatorname{span}\{e_0,e_1\}.
\]
Define the orthogonal map $W=\tfrac12 H$ using
\[
H=
\begin{pmatrix}
 1&-1&-1&-1\\
 1& 1& 1&-1\\
-1&-1& 1&-1\\
 1&-1& 1& 1
\end{pmatrix},
\qquad
V=W(U).
\]
Both principal angles between $U$ and $V$ are $\pi/4$, so the pair satisfies
the printed $\pi/3$ hypothesis.  The singular values of $I-W$ are
$(\sqrt2,\sqrt2,0,0)$, hence
\[
  \|I-W\|_1=2\sqrt2.
\]
For the direct rotation $R$ from $U$ to $V$, all four singular values of $I-R$
are $\sqrt{2-\sqrt2}$, giving
\[
  \|I-R\|_1=4\sqrt{2-\sqrt2}>2\sqrt2.
\]
Thus an admissible competitor has smaller trace-norm displacement than the
direct rotation.  The construction, its acuteness and principal-angle bounds,
the singular-value calculations, and the strict inequality are all checked in
Lean.

\subsection{Where the printed proof fails}

The counterexample also pinpoints the invalid step in the proof.  Write
$K=I-W$, and let $\Omega_1$ and $\Omega_2$ be the two principal-plane
projections used by Davis and Kahan.  Their equation~(4.3) requires, on this
configuration,
\begin{equation}
  \|K\|_4
  \geq
  \|K\Omega_1\|_2+\|K\Omega_2\|_2,
  \label{eq:dk43-false}
\end{equation}
where $\|\cdot\|_j$ denotes the Ky Fan $j$-norm.  The witness gives
\[
  \|K\|_4=2\sqrt2,
  \qquad
  \|K\Omega_1\|_2=\|K\Omega_2\|_2=2,
\]
so \cref{eq:dk43-false} would require $2\sqrt2\ge4$.  Orthogonality of the two
domain planes does not force the images under $K$ to be orthogonal, and the
singular-value sums of the restrictions therefore need not add.

Proposition~4.3 remains valid.  It concerns the squared displacement
$(I-W)^*(I-W)$, for which the formalization proves the source majorization
statement and its unitarily invariant norm consequences.  The two neighboring
claims therefore have different dispositions: squared-displacement
majorization is valid, while arbitrary unitarily invariant norm minimality of
the full displacement fails.

\subsection{A Q-norm repair}

The formal development proves a natural correction.  Call a unitarily invariant
norm $N$ a $Q$-norm if there is a unitarily invariant norm $M$ such that
\begin{equation}
  N(X)^2=M(X^*X)
  \label{eq:qnorm-def}
\end{equation}
for every $X$.  Proposition~4.3 applied to $M$ immediately yields the valid
full-displacement estimate
\begin{equation}
  N(I-R)\leq N(I-W)
  \label{eq:qnorm-repair}
\end{equation}
for the direct rotation $R$.  The Lean theorem proves
\cref{eq:qnorm-repair} over every \texttt{RCLike} scalar field.  Its angle
assumption is the acuteness condition needed to construct the direct rotation;
the printed $\pi/3$ threshold is absent.  Conversely, the trace norm used above is formally shown not
to satisfy this $Q$-norm predicate on the witness space.

Earlier work on canonical unitary completion already records failures for
Schatten norms below the quadratic regime, including trace-norm phenomena
\citep{bolotnikov2001minimal}.  The present counterexample is inside the
Davis--Kahan short-angle hypothesis and directly refutes Proposition~4.4 as
printed.  Later literature has repeated the short-angle claim
\citep{qiu2006grassmann}; the formal result therefore serves as a correction to
a still-propagated statement.

\section{Yu--Wang--Samworth as a statistical specialization}
\label{sec:yws}

YWS \citep{yu2015davis} gives a finite-dimensional population-gap version of
subspace perturbation theory.  Let
$\Sigma,\widehat\Sigma\in\mathbb{R}^{p\times p}$ be symmetric, with population
eigenvalues $\lambda_1\ge\cdots\ge\lambda_p$.  For
$1\le r\le s\le p$, write $d=s-r+1$ and let $V,\widehat V$ contain orthonormal
eigenvectors at indices $r,\ldots,s$.  With population gap
\[
  \Delta=
  \min\{\lambda_{r-1}-\lambda_r,\lambda_s-\lambda_{s+1}\}>0,
\]
YWS Theorem~2 gives
\begin{equation}
  \|\sin\Theta(\widehat V,V)\|_{\mathrm F}
  \le
  \frac{2\min\{\sqrt d\,\|\widehat\Sigma-\Sigma\|_{\mathrm{op}},
                    \|\widehat\Sigma-\Sigma\|_{\mathrm F}\}}
       {\Delta}.
  \label{eq:yws-sintheta}
\end{equation}
An aligned-frame bound follows with an additional factor $\sqrt2$.
Theorem~3 applies the same population-gap principle to left and right singular
subspaces through the corresponding Gram matrices.  The formalization checks
Theorems~1--3 and the supporting identities and conventions needed for those
claims; Appendix~\ref{app:yws-coverage} gives the exact scope.

Two printed defects appear in this comparison.  Equation~(4) omits a square.
For real unit vectors $u,v$, the corrected identity is
\begin{equation}
(2\langle v,u\rangle)^2(1-\langle v,u\rangle^2)
=
\frac14\|v-u\|^2(2-\|v-u\|^2)^2(4-\|v-u\|^2).
\label{eq:yws-equation4-corrected}
\end{equation}
Lean verifies both a numerical counterexample to the printed equation and the
corrected identity in general.

The second defect is the boundary convention in Theorem~3.  The assignment
$\sigma^2_{\operatorname{rank}(A)+1}:=-\infty$ can make the denominator
infinite when the selected block reaches the rank boundary, giving a zero upper
bound even when the compared singular subspaces differ.  The Gram-matrix proof
uses the ambient spectrum, where singular values after the rank and before the
ambient dimension are zero.  The corrected theorem extends the singular values
by zero through the ambient dimension and places the terminal convention after
that dimension.  A checked $2\times2$ rank-one example witnesses the failure
of the printed convention.

\section{AI-assisted formalization and the changing tooling landscape}
\label{sec:ai-landscape}

The starting point of this project differs substantially from the traditional
model of a formalization led by an expert in both the mathematics and the proof
assistant.  At the outset, the primary author had no prior experience with
formal theorem proving and had not studied most of the spectral and operator
theory needed for Davis--Kahan beyond basic linear algebra and set theory.  The
project nevertheless reached a checked formal treatment of the result-level
content of the 1970 paper, including its unbounded and infinite-dimensional
arguments.  The gap between that starting background and the resulting formal scope is an
observation about the practical capabilities of contemporary formalization
tools.  For the primary author, a project at this level of operator theory would
not have been attempted without AI assistance; the tools made it possible to
learn Lean and the required mathematics while constructing the development.

Large language models were used throughout the work.  ChatGPT, Claude Code, and
Codex-family agents explained unfamiliar mathematical material, proposed proof
decompositions, searched the repository and external formal libraries, wrote
Lean definitions and proofs, used compiler feedback to repair them, compared
formal statements with source passages, and helped draft and edit the
manuscript.  Formalization also became a way to learn the underlying
mathematics: definitions and domain conditions that could be passed over during
informal reading had to be made explicit before Lean would accept the intended
argument.

The experience also exposed limits of this workflow.  During development, an
agent could produce a valid theorem and describe the task as complete even when
the statement covered only bounded operators, only finite-dimensional spaces,
or a narrower spectral-gap configuration.  Successful compilation certifies
the encoded proposition; it does not identify which informal theorem was
encoded.  Repeated source comparison, including reviews performed in fresh
model contexts and review by the human authors, was therefore part of the
workflow.  Proposition~4.4 illustrates the converse situation: persistent
proof failure eventually led to mathematical analysis of the source claim and
to a counterexample.  The weaker statements obtained along the way were not
accepted under the original label.

Long-horizon Lean agents and persistent formalization environments are now
being developed explicitly around compiler feedback, library search, planning,
and theorem reuse \citep{zhang2026leanmarathon,mathcode2026,autolean2026}.
Human--AI interaction in advanced mathematical formalization is also becoming a
subject of study \citep{collins2026humanai,anthropic2026zeta}.  This project
provides one detailed case in that changing landscape.  It does not isolate the
causal contribution of model capability, Mathlib growth, local library
construction, workflow design, or human review; no controlled comparison among
those factors was performed.  It does show that the amount of prior expertise
required to initiate a deep formalization can be much lower than the expertise
embodied in the finished development, provided that kernel checking is paired
with sustained mathematical review of what has actually been stated.

\section{Related work}
\label{sec:related}

Classical perturbation theory for invariant and singular subspaces includes the
Davis--Kahan results, Wedin's singular-subspace bounds, and later matrix and
operator treatments
\citep{davis1970rotation,wedin1972perturbation,stewart1990matrix}.
Sylvester-equation methods provide closely related estimates for operator
angles and unbounded self-adjoint operators
\citep{grubisic2005sylvester,grubisic2010tan}.  YWS packages a
finite-dimensional version around a population eigengap for statistical use
\citep{yu2015davis}; related spectral perturbation bounds are now standard in
high-dimensional statistics \citep{cape2019two,chen2021spectral}.

The formal development builds on Lean~4 and Mathlib
\citep{demoura2021lean4,mathlib2020}.  Large source-driven formalizations such
as Flyspeck and the Liquid Tensor Experiment show how verification of a
specific body of mathematics can require substantial library construction
\citep{hales2017kepler,scholze2022liquid}.  Spectra is technically close to the
present work because it develops spectral and operator infrastructure in Lean;
we directly adapted and generalized parts of that development
\citep{bornemann2026spectra}.

AI-assisted theorem proving has progressed from statement translation and
sketch-guided proving to long-horizon repository-scale workflows
\citep{wu2022autoformalization,jiang2023dsp,zhang2026leanmarathon}.
FormalAlign, ShadowBench, and EconCSLib study or expose the separate problem of
agreement between an informal claim and its formal statement
\citep{lu2025formalalign,han2026shadowbench,garg2026econcs}.  The Davis--Kahan
experience gives this issue a concrete operator-theoretic form: the
half-infinite gap cases and the modulus relation in
\cref{eq:modulus-sine} are mathematical conditions for matching the source,
even after Lean has checked the proof.

\section{Conclusion}
\label{sec:conclusion}

The formalization recovers the operator-theoretic breadth of Davis and Kahan's
1970 paper in Lean: direct rotations, unitarily invariant norm extremality,
Sylvester separation, generalized sine and tangent estimates, branch selection,
and the unbounded single-angle theory.  The full result-level accounting has
one mathematical exception.  Proposition~4.4 is false under its published
hypotheses; the checked four-dimensional counterexample locates the failure in
the proof, while Proposition~4.3 supplies the squared-displacement theorem from
which a valid $Q$-norm repair follows.

The project also clarifies several formalization choices that are easy to hide
in a matrix-only restatement.  Approximation numbers carry singular-value
arguments into infinite-dimensional operator ideals.  Domain-carrying partial
maps make the unbounded hypotheses visible.  Form-bounded Sylvester gaps encode
both finite and half-infinite spectral separation.  The rectangular sine map
used in the Lean theorem agrees with Davis and Kahan's positive angle operator
at the level required by the norm inequality because its modulus is exactly
$\sin\Theta_0$.

The reusable layer now contains much of the mathematics needed to express these
statements, so the paper-specific proofs can depend on general operator theory.  The YWS development
shows how a smaller population-gap theorem can reuse the same spectral geometry
and also demonstrates the value of checking algebraic identities and endpoint
conventions.  The AI-assisted development process adds a separate observation:
a formalizer can now begin far outside both the subject-matter and Lean expert
communities and still build a substantial checked development, while the burden
of deciding whether the checked statement is the intended mathematics remains
a human mathematical responsibility supported by explicit review.

\section*{AI assistance disclosure}
ChatGPT, Claude, and Codex-family models were used throughout the formalization
and manuscript preparation.  Their uses included mathematical discussion,
literature and library search, proof decomposition, Lean code generation,
compiler-driven proof repair, source-to-statement review, and drafting and
editing manuscript prose.  The human authors selected the mathematical claims,
directed the formalization and review, checked source interpretations, and take
responsibility for the content of the paper.  Appendix~\ref{app:formalization-process}
gives additional workflow and model details.

\section*{Acknowledgments}

This material is based upon work supported by the Defense Advanced Research
Project Agency (DARPA) under Contract No. HR001125CE017. Any opinions, findings
and conclusions or recommendations expressed in this material are those of the
author(s) and do not necessarily reflect the views of the Defense Advanced
Research Project Agency (DARPA).

\bibliographystyle{unsrtnat}
\bibliography{references}

\clearpage
\appendix

\section{Reading the Lean signatures}
\label{app:lean-guide}
\label{app:reading-lean}

A Lean theorem declaration has the form
\[
  \texttt{theorem name (parameters) (hypotheses) : conclusion := proof}.
\]
Parameters introduce the mathematical objects under discussion.  A hypothesis
such as \texttt{(hA : IsSelfAdjoint A)} supplies the corresponding fact;
bundled predicates such as \texttt{IsTrialResidual} and
\texttt{IsExactSpectralDecomposition} collect related assumptions.  Braces mark
arguments normally inferred by Lean, while square brackets request typeclass
instances supplying ambient algebraic, analytic, or topological structure.  The
proposition after the final colon is checked by Lean's kernel
\citep{demoura2021lean4}.

In \Cref{lst:dk-sintheta-headline}, \texttt{RCLike} is the Mathlib abstraction
covering $\mathbb{R}$ and $\mathbb{C}$.  A type
\texttt{ContinuousLinearMap K F E} is a continuous linear map.  A type
\texttt{LinearPMap K E E} is a linear partial map whose
domain is a submodule; this is the representation used for possibly unbounded
operators.  \texttt{N.Mem T} states that the ideal gauge associated with $N$ is
finite on $T$, and \texttt{N.gaugeReal T} is the corresponding real value.  The
suffix \texttt{Real} refers to the codomain of the gauge and not to the scalar
field.

\section{Davis--Kahan result inventory}
\label{app:dk-inventory}

\Cref{tab:dk-full-inventory} records the maintained result-level denominator.
Every row is currently verified in the ordinary Lean build and has accepted
semantic/source-correspondence review. The disposition column distinguishes a
proof of the source claim from the checked refutation of Proposition~4.4.

\begin{longtable}{>{\raggedright\arraybackslash}p{0.16\linewidth}>{\raggedright\arraybackslash}p{0.55\linewidth}>{\raggedright\arraybackslash}p{0.18\linewidth}}
\caption{The \DKResultCount{} Davis--Kahan result-level targets.}
\label{tab:dk-full-inventory}\\
\toprule
ID & Result & Disposition \\
\midrule
\endfirsthead
\toprule
ID & Result & Disposition \\
\midrule
\endhead
\\path{S2-sin-theta} & Single-angle sine theorem & proved at source scope \\
\\path{S2-tan-theta} & Single-angle tangent theorem & proved at source scope \\
\\path{S2-sin-two-theta} & Double-angle sine theorem & proved at source scope \\
\\path{S2-tan-two-theta} & Double-angle tangent theorem & proved at source scope \\
\\path{DK-3.1-prop} & Acute direct rotation existence and uniqueness & proved at source scope \\
\\path{DK-3.2-prop} & Nonacute existence criterion & proved at source scope \\
\\path{DK-3.3-prop} & Principal square-root characterization & proved at source scope \\
\\path{DK-3.4-prop} & Square as a direct rotation & proved at source scope \\
\\path{DK-3.1-thm} & Classification of pairs of subspaces & proved at source scope \\
\\path{DK-3.1-cor} & Compact classification by angle eigenvalues & proved at source scope \\
\\path{DK-3.5-prop} & Angle commutation and eigenspace geometry & proved at source scope \\
\\path{DK-3.2-cor} & Reversal symmetry & proved at source scope \\
\\path{DK-4.1-prop} & Pointwise and singular-value extremality of the direct rotation & proved at source scope \\
\\path{DK-4.1-cor} & Unitarily invariant norm minimality of restricted displacement & proved at source scope \\
\\path{DK-4.2-prop} & Basis-angle square-sum extremality & proved at source scope \\
\\path{DK-4.3-prop} & Squared-displacement unitarily invariant norm minimality & proved at source scope \\
\\path{DK-4.4-prop} & Full-displacement claim of Proposition~4.4 & refuted as printed; repaired separately \\
\\path{DK-5.1-thm} & Banach-space Sylvester lower bound & proved at source scope \\
\\path{DK-5.2-thm} & Semibounded self-adjoint Sylvester theorem & proved at source scope \\
\\path{DK-5.1-lem} & Strong-cutoff convergence of singular values & proved at source scope \\
\\path{DK-6.1-lem} & Direct-sum norm comparison and converse & proved at source scope \\
\\path{DK-6.2-lem} & Reflection-pinch contraction & proved at source scope \\
\\path{DK-6.1-prop} & Sine proof, ambient limitation, and symmetric sine theorem & proved at source scope \\
\\path{DK-6.1-thm} & Generalized sine theorem & proved at source scope \\
\\path{DK-6.2-thm} & Pairwise-gap square-norm sine theorem & proved at source scope \\
\\path{DK-6.3-thm} & Generalized tangent theorem and supporting machinery & proved at source scope \\
\\path{DK-6.3-lem} & Finite-rank near-maximizer leakage estimate & proved at source scope \\
\\path{DK-8.1-thm} & Branch selection and spectral repulsion & proved at source scope \\
\\path{DK-8.2-thm} & Smallness selects the acute branch & proved at source scope \\
\bottomrule
\end{longtable}

The result table is keyed to source results.  One source result can require many
Lean declarations, while a single reusable Lean theorem can support more than
one source-specific wrapper.  The denominator therefore remains stable under
refactoring of implementation modules.

\section{Source-correspondence review}
\label{app:source-review}

A result is accepted only after two distinct checks.

\paragraph{Mechanical verification.}
The declaration must resolve in the ordinary project build, and every theorem
on which it depends must be accepted by Lean's kernel.  This establishes the
formal proposition exactly as encoded.

\paragraph{Source comparison.}
The encoded proposition is compared with the cited source location.  The review
checks the mathematical objects, hypotheses, conclusion, scalar field,
dimension, norm class, boundary conventions, and operator-domain assumptions.
For a nonliteral representation, the review also records the theorem or
mathematical argument that justifies the correspondence.  Equation
\cref{eq:modulus-sine} is the main example: the rectangular map in the Lean
sine-theta theorem represents the source's positive sine operator only after a
modulus/polar-decomposition argument.

This separation prevents theorem names from serving as evidence of fidelity.
For example, the proved declaration
\texttt{sinTheta\_unbounded\_intervalExterior\_characterizedWitness\_rclike}
remains a useful finite-gap theorem, but it is not the selected source endpoint
because it does not include the half-infinite gap cases.  Conversely,
Proposition~4.4 is complete in the inventory even though its printed assertion
has no proof: the completed formal treatment consists of the transcribed claim
as a proposition, a checked counterexample/refutation, a diagnosis of the
source proof, and repaired theorems.

\section{YWS source-audit scope}
\label{app:yws-coverage}

The YWS audit uses a different denominator because the journal paper does not
claim a line-by-line formalization of every mathematical statement in that
article.  The maintained census has \YWSCensusEntryCount{} rows:
\YWSCensusPaperSourceCount{} rows associated with printed source material and
\YWSCensusExplicitAdditionCount{} explicitly marked additions.  All
\YWSCensusVerifiedCount{} rows resolve to declarations in the normal build.

The source rows include the conclusions of Theorems~1--3 and the equations,
lemmas, examples, and boundary conventions needed to check those conclusions.
Several rows can share one correction.  The census has
\YWSCensusCorrectedCount{} rows whose implementation is classified as corrected,
but the manuscript reports two printed YWS defects: the missing square in
equation~(4) and the rank-boundary convention in Theorem~3.

\section{Formalization procedure and AI assistance}
\label{app:formalization-process}

The development combined ordinary Lean compilation with extensive LLM
assistance.  Work on a source theorem typically began by reading the statement
and surrounding proof, identifying standing assumptions, and locating the
relevant Lean interfaces.  The target was then decomposed into intermediate
mathematical obligations.  Existing Mathlib, Tau Ceti, local, and external Lean
results were searched before introducing new foundations.  Candidate
definitions and proofs were compiled with Lean and revised in response to
elaboration or proof failures.  A separate pass compared the accepted theorem
statement with the source, including its operator domains, scalar field,
dimension, norm class, constants, gap hypotheses, and endpoint conventions.

LLMs participated at every stage of this cycle.  They transcribed and explained
source material, proposed decompositions, searched code and literature,
implemented Lean declarations, repaired proofs from compiler feedback, and
performed adversarial source comparisons.  Human authors selected the target
scope, decided which source interpretations to accept, evaluated proposed
counterexamples and repairs, and reopened statements when later review exposed
a mismatch.  The stopping condition for a source result was therefore a
kernel-checked formal treatment together with an accepted mathematical
comparison to the source, or a checked refutation when the printed statement
was false.

The systems used during the project are summarized in
\Cref{tab:formalization-systems}.  Model assignments changed during the
project.  The table therefore records observed uses without assigning a fixed
role to any one model.

\begin{table}[H]
\centering
\small
\begin{tabularx}{\linewidth}{>{\raggedright\arraybackslash}p{0.20\linewidth}>{\raggedright\arraybackslash}p{0.28\linewidth}X}
\toprule
System & Models used & Observed uses \\
\midrule
ChatGPT chat interface & GPT-5.6 Sol & Mathematical discussion, proof planning, source comparison, review, literature search, and manuscript drafting and editing. \\
Claude Code & Claude Opus 4.8; Claude Opus 5; Claude Fable 5 & Repository search, Lean implementation, compilation, proof repair, and long-running formalization tasks. \\
Codex & GPT-5.6 Sol; GPT-5.6 Terra; GPT-5 Codex; GPT-5.6 Codex & Repository search, Lean implementation, compilation, repair, review, and proposed repository changes. \\
\bottomrule
\end{tabularx}
\caption{AI systems and model families used during the formalization and
manuscript preparation.  The prose of the paper was drafted and edited with AI
assistance; the authors reviewed and take responsibility for the submitted
text.}
\label{tab:formalization-systems}
\end{table}

Lean's kernel checks the proof terms produced by this process.  It does not
check that a theorem statement is the intended transcription of a paper.  The
statement-review procedure in Appendix~\ref{app:source-review} addresses that
second question.  Some early valid Lean statements covered narrower bounded,
finite-dimensional, or spectral-gap cases than the corresponding
Davis--Kahan theorem, so those statements were reopened after source review.

\section{Artifact organization}
\label{app:artifact}

The paper-facing Davis--Kahan declarations are under
\path{DavisKahan/Sources/DavisKahan1970}.  Reusable mathematical infrastructure
is primarily under \path{ForTauCeti}; its declarations already use final
\texttt{TauCeti.*} namespaces.  The maintained Davis--Kahan result inventory is
\path{dev/davis-kahan-1970-formalization-result-inventory.json}.  The YWS source
census is maintained separately under the YWS package and is summarized for
this manuscript by the generated census snapshot.  This organization permits
source-audit records and paper-specific counterexamples to remain local to the
paper package while reusable analysis is prepared for upstream integration.


\end{document}
